\documentclass[report, 12pt]{uftreport}

\usepackage{amsmath}

\title{Laboratório 5 -- Barrel Shifter e Somador Carry Look-Ahead}

\author{Pedro Lucas}{Moreira}

\advisor{Prof.}{Tiago}{Almeida}{}

\department{CC}
\class{Linguagens de Descrição de Hardware -- 2026/2}

\date{29}{9}{2026}

\keyword{VHDL}
\keyword{Barrel Shifter}
\keyword{Rotação Binária}
\keyword{Multiplexador}
\keyword{Carry Look-Ahead}
\keyword{Timing Analyzer}

\begin{document}

\frontmatter
\maketitle
\makefrontpage

\mainmatter
\ChapterStart{first}{Questão 1 -- Barrel Shifter de 4 Bits}

\chapter{Questão 1 -- Barrel Shifter de 4 Bits}

A rotação binária reposiciona os bits de uma palavra de forma circular, mantendo a ordem. Dada a palavra $W = w_3w_2w_1w_0$ e o valor de rotacionamento $S = s_1s_0$, a saída $Y = y_3y_2y_1y_0$ segue a Tabela~\ref{tab:barrel}.

\begin{table}[htb]
  \centering
  \caption{Tabela verdade do \textit{barrel shifter} de 4 bits.}
  \label{tab:barrel}
  \begin{tabular}{cc|cccc}
    \toprule
    $s_1$ & $s_0$ & $y_3$ & $y_2$ & $y_1$ & $y_0$ \\
    \midrule
    0 & 0 & $w_3$ & $w_2$ & $w_1$ & $w_0$ \\
    0 & 1 & $w_0$ & $w_3$ & $w_2$ & $w_1$ \\
    1 & 0 & $w_1$ & $w_0$ & $w_3$ & $w_2$ \\
    1 & 1 & $w_2$ & $w_1$ & $w_0$ & $w_3$ \\
    \bottomrule
  \end{tabular}
\end{table}

Observando cada coluna da tabela, cada bit de saída é sempre um bit da entrada deslocado de $S$ posições:

\[
  y_i = w_{(i + S) \bmod 4}, \qquad i \in \{0, 1, 2, 3\}
\]

\section{Item (a) -- Implementação em VHDL}

Seguindo a dica do enunciado, cada bit de saída foi implementado como um multiplexador 4:1, com \texttt{ROTATION\_IN} ($S$) como seletor e os quatro bits de \texttt{WORD\_IN} ($W$) como entradas de dados. Cada multiplexador é descrito por uma atribuição \texttt{with ... select}, cuja ordem das entradas é tirada diretamente da coluna correspondente da Tabela~\ref{tab:barrel}.

Arquivo \texttt{a.vhdl}:

\begin{lstlisting}[language=VHDL]
library ieee;
use ieee.std_logic_1164.all;

entity rotation_binary is
    port (
        WORD_IN     : in  std_logic_vector(3 downto 0);
        ROTATION_IN : in  std_logic_vector(1 downto 0);
        ROTATED_OUT : out std_logic_vector(3 downto 0)
    );
end entity rotation_binary;

architecture rtl of rotation_binary is
begin
    with ROTATION_IN select
        ROTATED_OUT(0) <= WORD_IN(0) when "00",
                          WORD_IN(1) when "01",
                          WORD_IN(2) when "10",
                          WORD_IN(3) when others;

    with ROTATION_IN select
        ROTATED_OUT(1) <= WORD_IN(1) when "00",
                          WORD_IN(2) when "01",
                          WORD_IN(3) when "10",
                          WORD_IN(0) when others;

    with ROTATION_IN select
        ROTATED_OUT(2) <= WORD_IN(2) when "00",
                          WORD_IN(3) when "01",
                          WORD_IN(0) when "10",
                          WORD_IN(1) when others;

    with ROTATION_IN select
        ROTATED_OUT(3) <= WORD_IN(3) when "00",
                          WORD_IN(0) when "01",
                          WORD_IN(1) when "10",
                          WORD_IN(2) when others;
end architecture rtl;
\end{lstlisting}

\subsection{Simulação}

A simulação mantém uma palavra fixa em \texttt{WORD\_IN} e percorre os quatro valores de \texttt{ROTATION\_IN}, com 10 ns entre cada estímulo. São usadas duas palavras: \texttt{1000}, com um único bit em 1 para deixar a rotação visível, e \texttt{1011}.

Arquivo \texttt{tb\_rotation\_binary.vhdl}, usado para a simulação:

\begin{lstlisting}[language=VHDL]
library ieee;
use ieee.std_logic_1164.all;
use ieee.numeric_std.all;

entity tb_rotation_binary is
end entity tb_rotation_binary;

architecture sim of tb_rotation_binary is
    signal WORD_IN     : std_logic_vector(3 downto 0);
    signal ROTATION_IN : std_logic_vector(1 downto 0);
    signal ROTATED_OUT : std_logic_vector(3 downto 0);
begin
    uut: entity work.rotation_binary
        port map (
            WORD_IN     => WORD_IN,
            ROTATION_IN => ROTATION_IN,
            ROTATED_OUT => ROTATED_OUT
        );

    stim: process
    begin
        -- palavra fixa com bits distintos para visualizar a rotacao
        WORD_IN <= "1000";
        for s in 0 to 3 loop
            ROTATION_IN <= std_logic_vector(to_unsigned(s, 2));
            wait for 10 ns;
        end loop;

        WORD_IN <= "1011";
        for s in 0 to 3 loop
            ROTATION_IN <= std_logic_vector(to_unsigned(s, 2));
            wait for 10 ns;
        end loop;
        wait;
    end process;
end architecture sim;
\end{lstlisting}

Resultados esperados, obtidos a partir da Tabela~\ref{tab:barrel}:

\begin{table}[htb]
  \centering
  \caption{Saídas esperadas na simulação.}
  \label{tab:sim}
  \begin{tabular}{c|c|c}
    \toprule
    \texttt{WORD\_IN} & \texttt{ROTATION\_IN} & \texttt{ROTATED\_OUT} \\
    \midrule
    1000 & 00 & 1000 \\
    1000 & 01 & 0100 \\
    1000 & 10 & 0010 \\
    1000 & 11 & 0001 \\
    \midrule
    1011 & 00 & 1011 \\
    1011 & 01 & 1101 \\
    1011 & 10 & 1110 \\
    1011 & 11 & 0111 \\
    \bottomrule
  \end{tabular}
\end{table}

% TODO: inserir a forma de onda da simulação
% \begin{figure}[htb]
%   \centering
%   \includegraphics[width=0.9\textwidth]{../01/simulacao.png}
%   \caption{Forma de onda da simulação do \textit{barrel shifter}.}
%   \label{fig:sim}
% \end{figure}

\section{Item (b) -- Execução na DE1 com o demo\_setup}

O \texttt{rotation\_binary} foi instanciado no \texttt{demo\_setup}. Os \textit{switches} \texttt{SW(3 downto 0)} fornecem a palavra $W$, \texttt{SW(5 downto 4)} fornecem o rotacionamento $S$, e a saída $Y$ é exibida nos LEDs vermelhos \texttt{LEDR3} a \texttt{LEDR0}.

\begin{table}[htb]
  \centering
  \caption{Mapeamento de entradas e saídas na DE1.}
  \label{tab:pinos}
  \begin{tabular}{lll}
    \toprule
    Sinal & Função & Pino \\
    \midrule
    \texttt{SW[0]} & $w_0$ & \texttt{PIN\_L22} \\
    \texttt{SW[1]} & $w_1$ & \texttt{PIN\_L21} \\
    \texttt{SW[2]} & $w_2$ & \texttt{PIN\_M22} \\
    \texttt{SW[3]} & $w_3$ & \texttt{PIN\_V12} \\
    \texttt{SW[4]} & $s_0$ & \texttt{PIN\_W12} \\
    \texttt{SW[5]} & $s_1$ & \texttt{PIN\_U12} \\
    \texttt{LEDR0} & $y_0$ & \texttt{PIN\_R20} \\
    \texttt{LEDR1} & $y_1$ & \texttt{PIN\_R19} \\
    \texttt{LEDR2} & $y_2$ & \texttt{PIN\_U19} \\
    \texttt{LEDR3} & $y_3$ & \texttt{PIN\_Y19} \\
    \bottomrule
  \end{tabular}
\end{table}

Arquivo \texttt{b.vhdl}:

\begin{lstlisting}[language=VHDL]
library ieee;
use ieee.std_logic_1164.all;

entity demo_setup is
    port (
        SW    : in  std_logic_vector(5 downto 0);
        LEDR0 : out std_logic;
        LEDR1 : out std_logic;
        LEDR2 : out std_logic;
        LEDR3 : out std_logic
    );
end entity demo_setup;

architecture rtl of demo_setup is
    signal rotated : std_logic_vector(3 downto 0);
begin
    shifter_inst: entity work.rotation_binary
        port map (
            WORD_IN     => SW(3 downto 0),
            ROTATION_IN => SW(5 downto 4),
            ROTATED_OUT => rotated
        );

    LEDR0 <= rotated(0);
    LEDR1 <= rotated(1);
    LEDR2 <= rotated(2);
    LEDR3 <= rotated(3);

end architecture rtl;
\end{lstlisting}

% TODO: inserir foto da placa em funcionamento
% \begin{figure}[htb]
%   \centering
%   \includegraphics[width=0.7\textwidth]{../01/placa.jpg}
%   \caption{\textit{Barrel shifter} em execução na DE1.}
%   \label{fig:placa}
% \end{figure}

\chapter{Questão 2 -- Somador Carry Look-Ahead}

\section{Item (a) -- CLA de 4 bits}

O somador \textit{carry look-ahead} (CLA) calcula todos os carries em paralelo, em vez de esperar o carry atravessar cada estágio como no \textit{ripple-carry}. Para cada bit $i$ são definidos os termos de geração e propagação:

\[
  G_i = A_i \cdot B_i, \qquad P_i = A_i \oplus B_i
\]

A partir deles, cada carry é expandido em função apenas de $G$, $P$ e $C_0 = C_{in}$:

\begin{align*}
  C_1 &= G_0 + P_0 C_0 \\
  C_2 &= G_1 + P_1 G_0 + P_1 P_0 C_0 \\
  C_3 &= G_2 + P_2 G_1 + P_2 P_1 G_0 + P_2 P_1 P_0 C_0 \\
  C_4 &= G_3 + P_3 G_2 + P_3 P_2 G_1 + P_3 P_2 P_1 G_0 + P_3 P_2 P_1 P_0 C_0
\end{align*}

A soma de cada bit é $S_i = P_i \oplus C_i$ e o carry-out é $C_{out} = C_4$.

Arquivo \texttt{a.vhdl}:

\begin{lstlisting}[language=VHDL]
library ieee;
use ieee.std_logic_1164.all;

entity cla_adder_4bit is
    port (
        A    : in  std_logic_vector(3 downto 0);
        B    : in  std_logic_vector(3 downto 0);
        Cin  : in  std_logic;
        S    : out std_logic_vector(3 downto 0);
        Cout : out std_logic
    );
end entity cla_adder_4bit;

architecture rtl of cla_adder_4bit is
    signal G : std_logic_vector(3 downto 0);
    signal P : std_logic_vector(3 downto 0);
    signal C : std_logic_vector(4 downto 0);
begin
    process (A, B)
    begin
        for i in 0 to 3 loop
            G(i) <= A(i) and B(i);
            P(i) <= A(i) xor B(i);
        end loop;
    end process;

    C(0) <= Cin;

    C(1) <= G(0) or (P(0) and C(0));

    C(2) <= G(1)
        or (P(1) and G(0))
        or (P(1) and P(0) and C(0));

    C(3) <= G(2)
        or (P(2) and G(1))
        or (P(2) and P(1) and G(0))
        or (P(2) and P(1) and P(0) and C(0));

    C(4) <= G(3)
        or (P(3) and G(2))
        or (P(3) and P(2) and G(1))
        or (P(3) and P(2) and P(1) and G(0))
        or (P(3) and P(2) and P(1) and P(0) and C(0));

    process (P, C)
    begin
        for i in 0 to 3 loop
            S(i) <= P(i) xor C(i);
        end loop;
    end process;

    Cout <= C(4);

end architecture rtl;
\end{lstlisting}

\subsection{Simulação}

A simulação é exaustiva: percorre as 512 combinações de \texttt{A}, \texttt{B} e \texttt{Cin}, com 10 ns entre cada estímulo, e compara \texttt{\{Cout, S\}} com o valor de $A + B + C_{in}$ calculado via \texttt{numeric\_std}. Foi executada via script \texttt{sim.do}, gerando o arquivo de onda \texttt{cla\_adder\_4bit.vcd}. O resultado foi \textbf{0 erros} nas 512 combinações.

Arquivo \texttt{tb\_cla\_adder\_4bit.vhdl}, usado para a simulação:

\begin{lstlisting}[language=VHDL]
library ieee;
use ieee.std_logic_1164.all;
use ieee.numeric_std.all;

-- Testbench exaustivo: percorre as 512 combinacoes de A, B e Cin
-- e compara a saida do CLA com a soma calculada via numeric_std.
entity tb_cla_adder_4bit is
end entity tb_cla_adder_4bit;

architecture sim of tb_cla_adder_4bit is
    signal A, B : std_logic_vector(3 downto 0);
    signal Cin  : std_logic;
    signal S    : std_logic_vector(3 downto 0);
    signal Cout : std_logic;
begin
    uut: entity work.cla_adder_4bit
        port map (
            A    => A,
            B    => B,
            Cin  => Cin,
            S    => S,
            Cout => Cout
        );

    stim: process
        variable expected : unsigned(4 downto 0);
        variable errors   : natural := 0;
    begin
        for c in 0 to 1 loop
            for a_val in 0 to 15 loop
                for b_val in 0 to 15 loop
                    A   <= std_logic_vector(to_unsigned(a_val, 4));
                    B   <= std_logic_vector(to_unsigned(b_val, 4));
                    if c = 1 then
                        Cin <= '1';
                    else
                        Cin <= '0';
                    end if;
                    wait for 10 ns;

                    expected := to_unsigned(a_val + b_val + c, 5);
                    if (Cout & S) /= std_logic_vector(expected) then
                        errors := errors + 1;
                        report "ERRO: A=" & integer'image(a_val)
                             & " B=" & integer'image(b_val)
                             & " Cin=" & integer'image(c)
                             severity error;
                    end if;
                end loop;
            end loop;
        end loop;

        report "Fim da simulacao: " & integer'image(errors) & " erro(s)"
            severity note;
        wait;
    end process;
end architecture sim;
\end{lstlisting}

\subsection{Comparação de tempo com o ripple-carry de 4 bits}

O CLA e o \texttt{ripple\_carry\_adder4} do laboratório anterior foram compilados nas mesmas condições: \textit{top-level} isolado, sem pinos fixos, no Cyclone II EP2C20F484C7. O atraso foi medido com o TimeQuest (\texttt{report\_datasheet}, tabela \textit{Propagation Delay}), e o atraso do circuito é o maior valor dessa tabela.

\begin{table}[htb]
  \centering
  \caption{Caminho crítico dos somadores de 4 bits.}
  \label{tab:tempos-4bits}
  \begin{tabular}{llcc}
    \toprule
    Somador & Caminho crítico & Atraso & LEs \\
    \midrule
    CLA 4 bits          & \texttt{B[2]} $\rightarrow$ \texttt{Cout}     & 13{,}129 ns & 15 \\
    Ripple-carry 4 bits & \texttt{y[2]} $\rightarrow$ \texttt{overflow} & 12{,}475 ns & 10 \\
    \bottomrule
  \end{tabular}
\end{table}

Em 4 bits, o \textit{ripple-carry} teve o menor tempo, 0{,}654 ns (cerca de 5\%) a menos que o CLA. Em ambos o caminho crítico termina no carry-out, como esperado.

O resultado se explica pela forma como o Quartus sintetiza o circuito. As portas não são implementadas como foram descritas: a síntese reduz as duas descrições a funções lógicas e as mapeia em LUTs de 4 entradas. Com apenas 9 entradas (\texttt{A}, \texttt{B} e \texttt{Cin}), os dois somadores viram redes de LUTs de profundidade parecida, e a diferença estrutural entre \textit{ripple} e \textit{lookahead} praticamente desaparece. A maior parte dos cerca de 12 ns é atraso fixo de buffers de entrada/saída e roteamento até os pinos, igual para os dois. Além disso, as equações de carry expandidas fazem o CLA ocupar mais elementos lógicos (15 contra 10), espalhando mais a lógica no chip.

A vantagem do CLA aparece quando o número de bits cresce: no \textit{ripple-carry} o atraso aumenta com $N$ (no laboratório anterior, 19{,}860 ns em 8 bits e 71{,}831 ns em 64 bits), enquanto no CLA o carry não depende da cadeia de estágios anteriores.

\subsection{Execução na DE1 com o demo\_setup}

O \texttt{cla\_adder\_4bit} foi instanciado no \texttt{demo\_setup}, com os \textit{switches} como entradas e os LEDs vermelhos como saídas.

\begin{table}[htb]
  \centering
  \caption{Mapeamento de entradas e saídas na DE1 (CLA).}
  \label{tab:pinos-cla}
  \begin{tabular}{lll}
    \toprule
    Sinal & Função & Pino \\
    \midrule
    \texttt{SW[3..0]} & $A_3..A_0$ & \texttt{V12}, \texttt{M22}, \texttt{L21}, \texttt{L22} \\
    \texttt{SW[7..4]} & $B_3..B_0$ & \texttt{M2}, \texttt{U11}, \texttt{U12}, \texttt{W12} \\
    \texttt{SW[8]}    & $C_{in}$   & \texttt{M1} \\
    \texttt{LEDR3..LEDR0} & $S_3..S_0$ & \texttt{Y19}, \texttt{U19}, \texttt{R19}, \texttt{R20} \\
    \texttt{LEDR4}    & $C_{out}$  & \texttt{T18} \\
    \bottomrule
  \end{tabular}
\end{table}

Arquivo \texttt{demo\_setup.vhd}:

\begin{lstlisting}[language=VHDL]
library ieee;
use ieee.std_logic_1164.all;

-- ITEM (a): instancia o cla_adder_4bit na placa DE1.
--   SW(3 downto 0) = entrada A
--   SW(7 downto 4) = entrada B
--   SW(8)          = Cin
--   LEDR3..LEDR0   = soma S
--   LEDR4          = Cout
entity demo_setup is
    port (
        SW    : in  std_logic_vector(8 downto 0);
        LEDR0 : out std_logic;
        LEDR1 : out std_logic;
        LEDR2 : out std_logic;
        LEDR3 : out std_logic;
        LEDR4 : out std_logic
    );
end entity demo_setup;

architecture rtl of demo_setup is
    signal sum_result : std_logic_vector(3 downto 0);
    signal carry_out  : std_logic;
begin
    adder_inst: entity work.cla_adder_4bit
        port map (
            A    => SW(3 downto 0),
            B    => SW(7 downto 4),
            Cin  => SW(8),
            S    => sum_result,
            Cout => carry_out
        );

    LEDR0 <= sum_result(0);
    LEDR1 <= sum_result(1);
    LEDR2 <= sum_result(2);
    LEDR3 <= sum_result(3);
    LEDR4 <= carry_out;

end architecture rtl;
\end{lstlisting}

\section{Item (b) -- CLA parcial de 8 bits}

O somador de 8 bits com CLA parcial é formado por dois blocos CLA de 4 bits ligados em cascata: o carry-out do bloco menos significativo, $C_4$, é o carry-in do bloco mais significativo. Dentro de cada bloco os carries são calculados em paralelo; o único elo em série é $C_4$. As equações do segundo bloco são as do primeiro deslocadas de 4 bits, com $C_4$ no lugar de $C_0$:

\begin{align*}
  C_5 &= G_4 + P_4 C_4 \\
  C_6 &= G_5 + P_5 G_4 + P_5 P_4 C_4 \\
  C_7 &= G_6 + P_6 G_5 + P_6 P_5 G_4 + P_6 P_5 P_4 C_4 \\
  C_8 &= G_7 + P_7 G_6 + P_7 P_6 G_5 + P_7 P_6 P_5 G_4 + P_7 P_6 P_5 P_4 C_4
\end{align*}

Os cálculos de $G$, $P$ e da soma usam \texttt{for ... generate}, para que cada bloco dirija apenas os seus bits dos vetores. A saída \texttt{Cout\_half\_1} expõe o carry entre os blocos ($C_4$) e \texttt{Cout\_half\_2} é o carry-out final ($C_8$).

Arquivo \texttt{b.vhdl}:

\begin{lstlisting}[language=VHDL]
library ieee;
use ieee.std_logic_1164.all;

entity cla_adder_8bit_with_4bit is
    port (
        A           : in  std_logic_vector(7 downto 0);
        B           : in  std_logic_vector(7 downto 0);
        Cin         : in  std_logic;
        S           : out std_logic_vector(7 downto 0);
        Cout_half_1 : out std_logic;
        Cout_half_2 : out std_logic
    );
end entity cla_adder_8bit_with_4bit;

architecture rtl of cla_adder_8bit_with_4bit is
    signal G : std_logic_vector(7 downto 0);
    signal P : std_logic_vector(7 downto 0);
    signal C : std_logic_vector(8 downto 0);
begin
    gp_low: for i in 0 to 3 generate
        G(i) <= A(i) and B(i);
        P(i) <= A(i) xor B(i);
    end generate gp_low;

    C(0) <= Cin;

    C(1) <= G(0) or (P(0) and C(0));

    C(2) <= G(1)
        or (P(1) and G(0))
        or (P(1) and P(0) and C(0));

    C(3) <= G(2)
        or (P(2) and G(1))
        or (P(2) and P(1) and G(0))
        or (P(2) and P(1) and P(0) and C(0));

    C(4) <= G(3)
        or (P(3) and G(2))
        or (P(3) and P(2) and G(1))
        or (P(3) and P(2) and P(1) and G(0))
        or (P(3) and P(2) and P(1) and P(0) and C(0));

    sum_low: for i in 0 to 3 generate
        S(i) <= P(i) xor C(i);
    end generate sum_low;

    Cout_half_1 <= C(4);

    gp_high: for i in 4 to 7 generate
        G(i) <= A(i) and B(i);
        P(i) <= A(i) xor B(i);
    end generate gp_high;

    C(5) <= G(4) or (P(4) and C(4));

    C(6) <= G(5)
        or (P(5) and G(4))
        or (P(5) and P(4) and C(4));

    C(7) <= G(6)
        or (P(6) and G(5))
        or (P(6) and P(5) and G(4))
        or (P(6) and P(5) and P(4) and C(4));

    C(8) <= G(7)
        or (P(7) and G(6))
        or (P(7) and P(6) and G(5))
        or (P(7) and P(6) and P(5) and G(4))
        or (P(7) and P(6) and P(5) and P(4) and C(4));

    sum_high: for i in 4 to 7 generate
        S(i) <= P(i) xor C(i);
    end generate sum_high;

    Cout_half_2 <= C(8);
end architecture rtl;
\end{lstlisting}

A implementação foi verificada com um testbench exaustivo, nos mesmos moldes do item (a), percorrendo as 131.072 combinações de \texttt{A}, \texttt{B} e \texttt{Cin}: 0 erros.

\section{Item (c) -- CLA puro de 8 bits}

No CLA puro não há nenhuma ligação em cascata: as equações de carry do item (a) são estendidas até $C_8$, e todo carry é escrito apenas em função de $G$, $P$ e $C_0$. De forma geral:

\[
  C_{k} = G_{k-1} + \sum_{j=0}^{k-2} \left( \prod_{m=j+1}^{k-1} P_m \right) G_j + \left( \prod_{m=0}^{k-1} P_m \right) C_0
\]

Por exemplo, o carry-out final é:

\[
  C_8 = G_7 + P_7 G_6 + P_7 P_6 G_5 + \dots + P_7 P_6 P_5 P_4 P_3 P_2 P_1 G_0 + P_7 P_6 P_5 P_4 P_3 P_2 P_1 P_0 C_0
\]

Arquivo \texttt{c.vhdl}:

\begin{lstlisting}[language=VHDL]
library ieee;
use ieee.std_logic_1164.all;

entity cla_adder_8bit is
    port (
        A    : in  std_logic_vector(7 downto 0);
        B    : in  std_logic_vector(7 downto 0);
        Cin  : in  std_logic;
        S    : out std_logic_vector(7 downto 0);
        Cout : out std_logic
    );
end entity cla_adder_8bit;

architecture rtl of cla_adder_8bit is
    signal G : std_logic_vector(7 downto 0);
    signal P : std_logic_vector(7 downto 0);
    signal C : std_logic_vector(8 downto 0);
begin
    process (A, B)
    begin
        for i in 0 to 7 loop
            G(i) <= A(i) and B(i);
            P(i) <= A(i) xor B(i);
        end loop;
    end process;

    C(0) <= Cin;

    C(1) <= G(0) or (P(0) and C(0));

    C(2) <= G(1)
        or (P(1) and G(0))
        or (P(1) and P(0) and C(0));

    C(3) <= G(2)
        or (P(2) and G(1))
        or (P(2) and P(1) and G(0))
        or (P(2) and P(1) and P(0) and C(0));

    C(4) <= G(3)
        or (P(3) and G(2))
        or (P(3) and P(2) and G(1))
        or (P(3) and P(2) and P(1) and G(0))
        or (P(3) and P(2) and P(1) and P(0) and C(0));

    C(5) <= G(4)
        or (P(4) and G(3))
        or (P(4) and P(3) and G(2))
        or (P(4) and P(3) and P(2) and G(1))
        or (P(4) and P(3) and P(2) and P(1) and G(0))
        or (P(4) and P(3) and P(2) and P(1) and P(0) and C(0));

    C(6) <= G(5)
        or (P(5) and G(4))
        or (P(5) and P(4) and G(3))
        or (P(5) and P(4) and P(3) and G(2))
        or (P(5) and P(4) and P(3) and P(2) and G(1))
        or (P(5) and P(4) and P(3) and P(2) and P(1) and G(0))
        or (P(5) and P(4) and P(3) and P(2) and P(1) and P(0) and C(0));

    C(7) <= G(6)
        or (P(6) and G(5))
        or (P(6) and P(5) and G(4))
        or (P(6) and P(5) and P(4) and G(3))
        or (P(6) and P(5) and P(4) and P(3) and G(2))
        or (P(6) and P(5) and P(4) and P(3) and P(2) and G(1))
        or (P(6) and P(5) and P(4) and P(3) and P(2) and P(1) and G(0))
        or (P(6) and P(5) and P(4) and P(3) and P(2) and P(1) and P(0) and C(0));

    C(8) <= G(7)
        or (P(7) and G(6))
        or (P(7) and P(6) and G(5))
        or (P(7) and P(6) and P(5) and G(4))
        or (P(7) and P(6) and P(5) and P(4) and G(3))
        or (P(7) and P(6) and P(5) and P(4) and P(3) and G(2))
        or (P(7) and P(6) and P(5) and P(4) and P(3) and P(2) and G(1))
        or (P(7) and P(6) and P(5) and P(4) and P(3) and P(2) and P(1) and G(0))
        or (P(7) and P(6) and P(5) and P(4) and P(3) and P(2) and P(1) and P(0) and C(0));

    process (P, C)
    begin
        for i in 0 to 7 loop
            S(i) <= P(i) xor C(i);
        end loop;
    end process;

    Cout <= C(8);

end architecture rtl;
\end{lstlisting}

Assim como no item (b), a implementação foi verificada com as 131.072 combinações de entrada: 0 erros.

\section{Análise de tempos dos somadores de 8 bits}

Os três somadores de 8 bits foram compilados nas mesmas condições do item (a) e medidos com o TimeQuest. O \textit{ripple-carry} de 8 bits do laboratório anterior foi recompilado e resultou em 19{,}860 ns, exatamente o valor medido naquele laboratório, o que confirma que as condições de medição são as mesmas.

\begin{table}[htb]
  \centering
  \caption{Caminho crítico dos somadores de 8 bits.}
  \label{tab:tempos-8bits}
  \begin{tabular}{llcc}
    \toprule
    Somador & Caminho crítico & Atraso & LEs \\
    \midrule
    Ripple-carry (lab. 4) & \texttt{x[0]} $\rightarrow$ \texttt{overflow}    & 19{,}860 ns & 18 \\
    CLA parcial (b)       & \texttt{A[0]} $\rightarrow$ \texttt{Cout\_half\_2} & 16{,}087 ns & 34 \\
    CLA puro (c)          & \texttt{B[2]} $\rightarrow$ \texttt{S[6]}        & 15{,}880 ns & 42 \\
    \bottomrule
  \end{tabular}
\end{table}

O menor tempo é o do \textbf{CLA puro}. Diferente do resultado em 4 bits, em 8 bits as duas versões de CLA são mais rápidas que o \textit{ripple-carry}:

\begin{itemize}
  \item \textbf{CLA parcial $\times$ ripple-carry:} 3{,}773 ns a menos (cerca de 19\%). No \textit{ripple-carry} o carry precisa atravessar 8 somadores completos em série; no CLA parcial ele só passa em série pelo elo $C_4$ entre os dois blocos. O caminho crítico ainda vai do bit 0 até o carry final, atravessando essa cascata.
  \item \textbf{CLA puro $\times$ CLA parcial:} 0{,}207 ns a menos (cerca de 1\%). Como nenhum carry depende de outro carry, o elo $C_4$ deixa de existir, e o caminho crítico passa a terminar em um bit da soma (\texttt{S[6]}) e não mais no carry-out. O ganho é pequeno porque a síntese já reorganiza as duas descrições em LUTs de 4 entradas, e $C_8$, com 17 entradas, precisa de vários níveis de LUT em qualquer das duas formas.
  \item \textbf{CLA puro $\times$ ripple-carry:} 3{,}980 ns a menos (cerca de 20\%).
\end{itemize}

A velocidade tem custo em área. O CLA puro ocupa 42 elementos lógicos, contra 34 do parcial e 18 do \textit{ripple-carry}, porque as equações de carry expandidas crescem com o número de bits ($C_8$ já tem 9 termos, o maior com 9 entradas). Por isso, para somadores maiores, costuma-se usar o CLA em blocos, como no item (b), ou de forma hierárquica, em vez de expandir todos os carries.

A análise completa dos tempos dos itens (a), (b) e (c) também está no arquivo \texttt{tempos.txt}.

\backmatter

\end{document}
